% ---------------------------------------------------------------------------
% Geruest fuer eine neue Runde. Vor \end{document} einfuegen.
% Platzhalter: <NN> Rundennummer, <TITEL>, <SLUG> kebab-case des Titels.
% Jede eckige Klammer ist zu ersetzen oder zu loeschen -- nichts stehen lassen.
% ---------------------------------------------------------------------------

\clearpage
\section*{v4.7.0-r<NN> <TITEL>}
\addcontentsline{toc}{section}{v4.7.0-r<NN> <TITEL>}
\label{sec:r<NN>-<SLUG>}

\paragraph{Purpose.}
<Ein bis drei Saetze: was diese Runde tut und woran sie anschliesst.>

\noindent\textbf{Source-only status.}
This r<NN> layer is source-only.  It does not build a PDF, does not assign a new DOI,
and does not authorize a public release.  The intended reserved v4.7.0 Version DOI remains
\href{https://doi.org/10.5281/zenodo.20232911}{10.5281/zenodo.20232911}; the previous
baseline DOI \href{https://doi.org/10.5281/zenodo.20117361}{10.5281/zenodo.20117361}
remains provenance only.

% --------------------------------------------------------------- Hauptinhalt
\subsection*{r<NN> <ERSTER UNTERABSCHNITT>}
\label{subsec:r<NN>-<SLUG-1>}

<Aufbau des Objekts, Annahmen, Notation.>

\paragraph{Theoremlet r<NN>.}
<Aussage. Nur effektiv/lokal/frame-typisiert formulieren.
 Niemals: C5 = 0, 0_termdata => 0_all, C6/T2, W_gamma^glob, W_all.>

\paragraph{Proof sketch.}
<Beweisskizze mit Verweisen auf die Runden, aus denen die Voraussetzungen stammen.>
\(\square\)

% ------------------------------------------------------------- Guard-Bilanz
\subsection*{r<NN> guard update}
\label{subsec:r<NN>-guard-update}

The r<NN> layer does not spend any global guard:
\[
R_{\rm <BEREICH>}^{r<NN>}=(0,0,0),
\qquad
R_{\rm glob}^{(r<NN>,+)}=(0,1,1,1,1),
\]
with
\[
\Gamma_{\rm guard}^{r<NN>}=(1,1,1,1,1).
\]
The result does not imply exact finite \(C5\) zero closure, exact all-readout terminality,
\(C6/T2\), a global photon-clock/RH theorem, or global all-readout closure.

% Falls diese Runde doch einen Guard veraendert: den geaenderten Vektor hier
% eintragen UND begruenden, warum die Aenderung zulaessig ist. Eine solche
% Aenderung ist ein Ereignis und braucht die ausdrueckliche Zustimmung des Autors.
